% Biên dịch bằng XeLaTeX:  xelatex nxcred-crypto.tex
\documentclass[a4paper,11pt]{article}

\usepackage{fontspec}
\setmainfont{TeX Gyre Termes}
\setsansfont{TeX Gyre Heros}
\setmonofont{TeX Gyre Cursor}

\usepackage[margin=2.6cm]{geometry}
\usepackage{amsmath,amssymb,amsthm}
\usepackage{booktabs}
\usepackage{enumitem}
\usepackage{xcolor}
\usepackage[hidelinks]{hyperref}

\definecolor{ink}{HTML}{16171A}
\definecolor{muted}{HTML}{6E7079}
\definecolor{accent}{HTML}{2F5FE0}
\hypersetup{colorlinks=true,linkcolor=accent,urlcolor=accent}

\theoremstyle{definition}
\newtheorem{definition}{Định nghĩa}[section]
\theoremstyle{plain}
\newtheorem{theorem}[definition]{Định lý}
\newtheorem{lemma}[definition]{Bổ đề}
\newtheorem{proposition}[definition]{Mệnh đề}
\theoremstyle{remark}
\newtheorem*{remark}{Nhận xét}
\renewcommand{\proofname}{Chứng minh}

\newcommand{\Zn}{\mathbb{Z}_n^{*}}
\newcommand{\QR}{\mathrm{QR}_n}
\newcommand{\Hid}{\mathcal{H}}
\newcommand{\Rev}{\mathcal{R}}
\newcommand{\Hash}{\mathsf{H}}
\newcommand{\enc}{\mathsf{enc}}

\title{\sffamily\bfseries Lõi mật mã NX Cred\\[6pt]
\large\mdseries Chữ ký Camenisch--Lysyanskaya và bằng chứng không tiết lộ}
\author{}
\date{\small\color{muted}\today}

\begin{document}
\maketitle

\begin{abstract}
\noindent Tài liệu này trình bày đầy đủ phần toán học của NX Cred: cấu trúc nhóm, lược đồ
cam kết, giao thức sigma dùng khi xin cấp thẻ, chữ ký mù Camenisch--Lysyanskaya, và bằng
chứng trình bày có tiết lộ chọn lọc. Mọi công thức khớp một--một với mã nguồn trong
\texttt{issuer/issuer.py}, \texttt{wallet/wallet.py}, \texttt{verifier/verifier.py} và
\texttt{common.py}; các hằng số độ dài bit lấy đúng từ mã. Mỗi phép kiểm đều đi kèm chứng
minh tính đúng đắn.
\end{abstract}

\tableofcontents
\bigskip\hrule\bigskip

\section{Ký hiệu và tham số}

Mọi phép nhân và luỹ thừa đều hiểu là trong $\Zn$, tức modulo $n$, trừ khi ghi rõ là phép
toán trên $\mathbb{Z}$.

\begin{center}
\begin{tabular}{@{}llr@{}}
\toprule
\textbf{Ký hiệu} & \textbf{Ý nghĩa} & \textbf{Độ dài} \\
\midrule
$p,q$ & hai safe prime, khoá riêng của issuer & $1024$ bit \\
$n=pq$ & modulus công khai & $\approx 2048$ bit \\
$\varphi=(p-1)(q-1)$ & chỉ issuer tính được & \\
$S$ & cơ số của phần ngẫu nhiên $v$ & \\
$R$ & cơ số của link secret & \\
$Z$ & đích của đẳng thức chữ ký & \\
$R_i$ & cơ số riêng của thuộc tính thứ $i$ & $8$ thuộc tính \\
$\mathit{ls}$ & link secret, chỉ ví biết & $256$ bit \\
$m_i$ & thuộc tính đã mã hoá thành số nguyên & \\
$(A,e,v)$ & credential: chữ ký CL trên bộ thuộc tính & \\
\bottomrule
\end{tabular}
\end{center}

\noindent Hằng số dùng khi trình bày:
\[
\ell_{e_0}=596,\quad
\ell_{\tilde e}=456,\quad
\ell_{\tilde v}=3060,\quad
\ell_{\tilde m}=593,\quad
\ell_r=2128 .
\]
Giai đoạn xin cấp dùng bộ số che riêng: $3488$ bit cho $\tilde v$ và $593$ bit cho
$\widetilde{\mathit{ls}}$.

\begin{remark}
Mọi số che đều dài hơn số bị che ít nhất $80$ bit. Đây là biên thống kê: phần dư ngẫu
nhiên nhấn chìm phần mang thông tin, xem Bổ đề~\ref{lem:stat}.
\end{remark}

\subsection{Mã hoá thuộc tính}

\begin{definition}[\texttt{common.encode\_attribute}]
Với giá trị $x$,
\[
\enc(x)=
\begin{cases}
\operatorname{int}(x), & \text{nếu } x \text{ là chuỗi toàn chữ số ASCII},\\[2pt]
\operatorname{int}\bigl(\mathrm{SHA\text{-}256}(x)\bigr), & \text{ngược lại.}
\end{cases}
\]
\end{definition}

\noindent Số căn cước giữ nguyên giá trị; họ tên, địa chỉ, ngày sinh băm thành số $256$
bit. Mọi phép toán phía sau chỉ làm việc với $m_i=\enc(x_i)$.

\subsection{Hàm thách thức}

\begin{definition}[\texttt{common.hash\_three}]
\[
\Hash(x,y,\mathit{nonce})=\operatorname{int}\bigl(\mathrm{SHA\text{-}256}(x \,\|\, y \,\|\, \mathit{nonce})\bigr).
\]
\end{definition}

\noindent Đây là biến đổi Fiat--Shamir: thay vì bên kiểm gửi thách thức rồi chờ trả lời,
bên chứng minh tự băm dữ liệu của chính mình ra thách thức. Vì không đoán trước được giá
trị băm nên không dựng sẵn được câu trả lời giả.

\section{Sinh khoá}

\begin{definition}[\texttt{issuer.setup}]
Chọn $p=2p'+1$ và $q=2q'+1$ với $p',q'$ cũng nguyên tố, mỗi số $1024$ bit. Đặt $n=pq$.
Lấy ngẫu nhiên $a,b,z,r_i \in [0,n)$ và đặt
\[
S=a^2,\quad R=b^2,\quad Z=z^2,\quad R_i=r_i^2 \pmod n .
\]
Công khai $(n,S,R,Z,\{R_i\})$; giữ kín $(p,q)$.
\end{definition}

\begin{proposition}[Vì sao bình phương mọi cơ số]
Mọi cơ số công khai đều nằm trong nhóm $\QR$ các thặng dư bình phương modulo $n$.
\end{proposition}

\begin{proof}
Theo định nghĩa, $S=a^2$ với $a\in\Zn$, nên $S\in\QR$; tương tự cho $R,Z,R_i$. Vì $\QR$
đóng với phép nhân và phép nghịch đảo, mọi tích và luỹ thừa của chúng cũng nằm trong
$\QR$.
\end{proof}

\begin{remark}
Tính chất này quan trọng vì nó bảo đảm mọi phần tử tham gia phép kiểm nằm cùng một nhóm
con, không có phần tử nào rơi ra ngoài làm hỏng đẳng thức. Việc chọn safe prime khiến
$|\QR|=p'q'$ có thừa số nguyên tố lớn, chặn các thuật toán khai thác nhóm bậc trơn
(Pohlig--Hellman).
\end{remark}

\section{Cam kết link secret}

\begin{definition}[\texttt{wallet.compute\_commitment}]
Ví sinh $\mathit{ls}\leftarrow\{0,1\}^{256}$ một lần và giữ mãi, sinh
$v'\leftarrow\{0,1\}^{2048}$ cho mỗi lần xin cấp, rồi tính
\[
u = S^{v'}\, R^{\mathit{ls}} \bmod n .
\]
\end{definition}

\begin{theorem}[Cam kết giấu kín]
\label{thm:hiding}
Giả sử $S$ sinh ra nhóm con chứa $R$, tức tồn tại $\alpha$ với $R=S^{\alpha}$. Khi đó với
mỗi giá trị giả định $\widehat{\mathit{ls}}$ đều tồn tại đúng một $\widehat{v}$ sao cho
$S^{\widehat v}R^{\widehat{\mathit{ls}}}=u$. Do đó $u$ không mang thông tin nào về
$\mathit{ls}$.
\end{theorem}

\begin{proof}
Đặt $k$ là bậc của $S$. Từ $R=S^{\alpha}$ ta có
$u=S^{v'+\alpha\,\mathit{ls}}$. Với $\widehat{\mathit{ls}}$ bất kỳ, chọn
\[
\widehat v \equiv v' + \alpha\,\mathit{ls} - \alpha\,\widehat{\mathit{ls}} \pmod{k}.
\]
Khi đó $S^{\widehat v}R^{\widehat{\mathit{ls}}}
= S^{v'+\alpha \mathit{ls}-\alpha\widehat{\mathit{ls}}}S^{\alpha\widehat{\mathit{ls}}}
= S^{v'+\alpha\mathit{ls}} = u$. Nghiệm là duy nhất modulo $k$ vì $S$ có bậc $k$. Vậy mọi
$\widehat{\mathit{ls}}$ đều tương thích với $u$ qua đúng một $\widehat v$, nên phân phối
của $u$ độc lập với $\mathit{ls}$.
\end{proof}

\begin{remark}
Đây là điểm mấu chốt của toàn bộ mô hình: \emph{issuer ký lên một thứ nó không đọc được}.
\end{remark}

\section{Giao thức sigma khi xin cấp}

Chỉ gửi $u$ là chưa đủ: kẻ gian có thể lấy $u$ của người khác rồi tự nhận. Ví phải chứng
minh mình biết cặp $(v',\mathit{ls})$ bên trong $u$ mà không nói ra cặp đó.

\begin{definition}[Bên chứng minh]
Ví lấy $\tilde v\leftarrow\{0,1\}^{3488}$,
$\widetilde{\mathit{ls}}\leftarrow\{0,1\}^{593}$ rồi tính
\[
\tilde u = S^{\tilde v} R^{\widetilde{\mathit{ls}}},\qquad
c=\Hash(u,\tilde u,\mathit{nonce}),
\]
\[
\hat v = \tilde v + c\,v' ,\qquad
\widehat{\mathit{ls}} = \widetilde{\mathit{ls}} + c\,\mathit{ls}
\qquad\text{(cộng trên }\mathbb{Z}\text{, không lấy modulo).}
\]
Gửi $(\mathit{nonce},u,c,\hat v,\widehat{\mathit{ls}})$. Không gửi $\tilde u$.
\end{definition}

\begin{definition}[Bên kiểm --- \texttt{issuer.verify\_proof}]
Issuer kiểm $1<u<n$ và $\gcd(u,n)=1$, rồi tính
\[
\tilde u' = S^{\hat v} R^{\widehat{\mathit{ls}}}\, u^{-c} \bmod n
\]
và chấp nhận khi $\Hash(u,\tilde u',\mathit{nonce})=c$.
\end{definition}

\begin{theorem}[Tính đầy đủ]
Nếu ví trung thực thì $\tilde u' = \tilde u$, do đó phép kiểm luôn thành công.
\end{theorem}

\begin{proof}
Thay trực tiếp các đáp số:
\begin{align*}
S^{\hat v} R^{\widehat{\mathit{ls}}} u^{-c}
&= S^{\tilde v + c v'}\; R^{\widetilde{\mathit{ls}} + c\,\mathit{ls}}\;
   \bigl(S^{v'}R^{\mathit{ls}}\bigr)^{-c} \\
&= S^{\tilde v} S^{c v'} R^{\widetilde{\mathit{ls}}} R^{c\,\mathit{ls}}
   S^{-c v'} R^{-c\,\mathit{ls}} \\
&= S^{\tilde v} R^{\widetilde{\mathit{ls}}} \;=\; \tilde u .
\end{align*}
Mọi thành phần nhân với $c$ triệt tiêu. Issuer dựng lại được $\tilde u$ mà chưa bao giờ
nhìn thấy nó, nên băm ra đúng $c$.
\end{proof}

\begin{theorem}[Tính chắc chắn --- trích xuất đặc biệt]
\label{thm:soundness}
Giả sử một bên chứng minh trả lời thành công hai thách thức khác nhau $c_1\neq c_2$ trên
cùng một $\tilde u$. Khi đó trích xuất được $\mathit{ls}$.
\end{theorem}

\begin{proof}
Từ hai lần trả lời hợp lệ:
\[
S^{\hat v_1} R^{\widehat{\mathit{ls}}_1} u^{-c_1} = \tilde u
= S^{\hat v_2} R^{\widehat{\mathit{ls}}_2} u^{-c_2}.
\]
Chuyển vế:
\[
S^{\hat v_1-\hat v_2}\,R^{\widehat{\mathit{ls}}_1-\widehat{\mathit{ls}}_2}
= u^{\,c_1-c_2}.
\]
Đặt $\Delta c=c_1-c_2\neq 0$. Nếu $\Delta c$ chia hết
$\widehat{\mathit{ls}}_1-\widehat{\mathit{ls}}_2$ và $\hat v_1-\hat v_2$ thì
\[
\mathit{ls}=\frac{\widehat{\mathit{ls}}_1-\widehat{\mathit{ls}}_2}{\Delta c},
\qquad
v'=\frac{\hat v_1-\hat v_2}{\Delta c}
\]
là một mở cam kết hợp lệ của $u$. Vậy ai trả lời được cho hai thách thức khác nhau thì
thực sự biết $\mathit{ls}$; ngược lại, kẻ không biết $\mathit{ls}$ chỉ có thể thành công
với xác suất bằng xác suất đoán trúng $c$.
\end{proof}

\begin{lemma}[Giấu kín thống kê]
\label{lem:stat}
Cho bí mật $s$ dài tối đa $\ell_s$ bit, thách thức $c$ dài tối đa $\ell_c$ bit, và số che
$\tilde s$ lấy đều trên $\{0,1\}^{\ell_{\tilde s}}$ với
$\ell_{\tilde s}\ge \ell_s+\ell_c+\kappa$. Khi đó khoảng cách thống kê giữa
$\tilde s + c\,s$ và một số ngẫu nhiên đều trên $\{0,1\}^{\ell_{\tilde s}}$ nhỏ hơn
$2^{-\kappa}$.
\end{lemma}

\begin{proof}
$\tilde s$ đều trên một khoảng độ dài $2^{\ell_{\tilde s}}$, còn $c\,s$ là một hằng số
nằm trong $[0,2^{\ell_s+\ell_c})$. Tổng là phân phối đều trên khoảng bị dịch đi. Hai
khoảng đều độ dài $2^{\ell_{\tilde s}}$ lệch nhau không quá $2^{\ell_s+\ell_c}$ nên phần
không giao nhau chiếm tỉ lệ tối đa $2^{\ell_s+\ell_c}/2^{\ell_{\tilde s}} \le 2^{-\kappa}$.
Đó chính là khoảng cách thống kê.
\end{proof}

\noindent Với $\kappa=80$ và các hằng số ở mục~1, mọi đáp số gửi đi đều thoả bổ đề này.

\section{Chữ ký mù}

\begin{definition}[\texttt{issuer.sign\_blindly}]
Issuer chọn số nguyên tố $e \in [2^{596},\,2^{596}+2^{119})$ với $\gcd(e,\varphi)=1$, lấy
$v''$ dài $2724$ bit với bit cao nhất bật, tính $d=e^{-1}\bmod\varphi$ rồi đặt
\[
Q = Z \cdot \Bigl( u\, S^{v''} \prod_i R_i^{m_i} \Bigr)^{-1} \bmod n,
\qquad
A = Q^{\,d} \bmod n .
\]
Trả về $(A,e,v'')$.
\end{definition}

\begin{remark}
Vế trong ngoặc gom đủ ba nhóm: cam kết $u$ của ví (chứa link secret), phần ngẫu nhiên
$v''$ của issuer, và toàn bộ thuộc tính. $A$ là căn bậc $e$ của $Q$, chỉ lấy được nhờ
$d$, tức nhờ biết $\varphi$, tức nhờ biết $p,q$.
\end{remark}

\section{Gỡ mù và đẳng thức chữ ký}

\begin{definition}[\texttt{wallet.unblind\_signature}]
Ví đặt $v=v'+v''$ và lưu credential $(A,e,v)$.
\end{definition}

\begin{theorem}[Đẳng thức chữ ký CL]
\label{thm:cl}
Với credential hợp lệ,
\[
Z \;=\; A^{e}\, S^{v}\, R^{\mathit{ls}} \prod_i R_i^{m_i} \pmod n .
\]
\end{theorem}

\begin{proof}
Theo cách dựng, $A=Q^{d}$ với $d e \equiv 1 \pmod{\varphi}$, nên theo định lý Euler
$A^{e}=Q^{de}=Q$. Thay $Q$ và $u=S^{v'}R^{\mathit{ls}}$:
\begin{align*}
A^{e} S^{v} R^{\mathit{ls}} \prod_i R_i^{m_i}
&= Z\Bigl(u\,S^{v''}\prod_i R_i^{m_i}\Bigr)^{-1} S^{v'+v''} R^{\mathit{ls}} \prod_i R_i^{m_i}\\
&= Z\Bigl(S^{v'}R^{\mathit{ls}}S^{v''}\prod_i R_i^{m_i}\Bigr)^{-1}
   S^{v'+v''} R^{\mathit{ls}} \prod_i R_i^{m_i}\\
&= Z .\qedhere
\end{align*}
\end{proof}

\noindent Đây là bất biến trung tâm của toàn hệ thống. Sai một thuộc tính, sai link
secret, hay chữ ký giả --- đều không cho ra $Z$. Ví kiểm đẳng thức này trước khi lưu
(\texttt{wallet.verify\_credential}) để issuer không thể trả về chữ ký hỏng rồi đổ lỗi
cho ví về sau.

\section{Bằng chứng trình bày}

Gọi $\Hid$ là tập thuộc tính giữ kín và $\Rev$ là tập tiết lộ, $\Hid \cup \Rev$ là toàn
bộ tám thuộc tính.

\subsection{Ngẫu nhiên hoá chữ ký}

\begin{definition}
Ví lấy $r\leftarrow\{0,1\}^{2128}$ rồi đặt
\[
A' = A\,S^{r},\qquad
v^{*} = v - e r \ \ (\text{trên }\mathbb{Z},\ \text{có thể âm}),\qquad
e^{*} = e - 2^{596}.
\]
\end{definition}

\begin{lemma}[Đẳng thức được bảo toàn]
\label{lem:rerand}
\[
Z = A'^{\,e}\, S^{v^{*}}\, R^{\mathit{ls}} \prod_i R_i^{m_i}.
\]
\end{lemma}

\begin{proof}
Từ $A'=A S^{r}$ suy ra $A = A' S^{-r}$, do đó $A^{e} = A'^{\,e} S^{-er}$. Thay vào
Định lý~\ref{thm:cl}:
\[
Z = A'^{\,e} S^{-er} S^{v} R^{\mathit{ls}} \prod_i R_i^{m_i}
  = A'^{\,e} S^{\,v-er} R^{\mathit{ls}} \prod_i R_i^{m_i}
  = A'^{\,e} S^{v^{*}} R^{\mathit{ls}} \prod_i R_i^{m_i}. \qedhere
\]
\end{proof}

\begin{remark}
$A'$ là thành phần duy nhất của credential rời khỏi ví, và nó khác nhau ở \emph{mỗi} lần
trình bày. Đây chính là tính không liên kết được: hai verifier gộp dữ liệu cũng không tìm
ra điểm chung. Ngoài ra $e^{*}<2^{119}$ nên số che $456$ bit đủ nhấn chìm nó theo
Bổ đề~\ref{lem:stat}.
\end{remark}

\subsection{Che và đáp số}

\begin{definition}
Lấy $\tilde e\leftarrow\{0,1\}^{456}$, $\tilde v\leftarrow\{0,1\}^{3060}$,
$\tilde m_{\mathit{ls}}\leftarrow\{0,1\}^{593}$ và
$\tilde m_i\leftarrow\{0,1\}^{593}$ với mỗi $i\in\Hid$. Đặt
\[
T = A'^{\,\tilde e}\, S^{\tilde v}\, R^{\tilde m_{\mathit{ls}}}
    \prod_{i\in\Hid} R_i^{\tilde m_i},
\qquad
c = \Hash(T, A', n_v),
\]
\[
\hat e = \tilde e + c\,e^{*},\quad
\hat v = \tilde v + c\,v^{*},\quad
\hat m_{\mathit{ls}} = \tilde m_{\mathit{ls}} + c\,\mathit{ls},\quad
\hat m_i = \tilde m_i + c\,m_i \ (i\in\Hid).
\]
Gửi cho verifier: $A',c,\hat e,\hat v,\hat m_{\mathit{ls}},\{\hat m_i\}_{i\in\Hid}$ và
các cặp (giá trị thô, $m_i$) với $i\in\Rev$.
\end{definition}

\section{Phép kiểm của verifier}

Verifier lấy $(n,S,R,Z,\{R_i\})$ \emph{từ chain}, không nhận từ bên trình bày.

\subsection{Kiểm biên}

\[
1<A'<n,\qquad 0\le \hat e < 2^{457},\qquad 0\le c<2^{256},
\]
tập tiết lộ khớp đúng yêu cầu đã phát, tập ẩn khớp phần còn lại của schema, và với mỗi
$i\in\Rev$ phải có $\enc(\text{giá trị thô})=m_i$. Phép kiểm cuối buộc giá trị chữ mà
verifier đọc phải khớp với số đưa vào phép tính.

\subsection{Dựng lại đích và khôi phục}

\begin{definition}
\[
D = Z \cdot \Bigl(\prod_{i\in\Rev} R_i^{m_i}\Bigr)^{-1} A'^{\,-2^{596}} \bmod n,
\qquad
\widehat T = D^{-c}\, A'^{\,\hat e}\, S^{\hat v}\, R^{\hat m_{\mathit{ls}}}
             \prod_{i\in\Hid} R_i^{\hat m_i}.
\]
Chấp nhận khi $\Hash(\widehat T, A', n_v) = c$.
\end{definition}

\begin{lemma}[$D$ đúng bằng phần bí mật]
\label{lem:D}
\[
D = A'^{\,e^{*}}\, S^{v^{*}}\, R^{\mathit{ls}} \prod_{i\in\Hid} R_i^{m_i}.
\]
\end{lemma}

\begin{proof}
Từ Bổ đề~\ref{lem:rerand} và $e=e^{*}+2^{596}$:
\[
Z = A'^{\,e^{*}+2^{596}} S^{v^{*}} R^{\mathit{ls}}
    \prod_{i\in\Hid} R_i^{m_i} \prod_{i\in\Rev} R_i^{m_i}.
\]
Chuyển hai thừa số mà verifier tự tính được sang vế trái:
\[
Z \Bigl(\prod_{i\in\Rev} R_i^{m_i}\Bigr)^{-1} A'^{\,-2^{596}}
= A'^{\,e^{*}} S^{v^{*}} R^{\mathit{ls}} \prod_{i\in\Hid} R_i^{m_i}.
\]
Vế trái đúng bằng $D$.
\end{proof}

\begin{theorem}[Tính đúng đắn của phép kiểm]
Với bằng chứng trung thực, $\widehat T = T$.
\end{theorem}

\begin{proof}
Thay các đáp số và tách theo luỹ thừa của $c$:
\begin{align*}
A'^{\,\hat e} S^{\hat v} R^{\hat m_{\mathit{ls}}} \prod_{i\in\Hid} R_i^{\hat m_i}
&= A'^{\,\tilde e + c e^{*}} S^{\tilde v + c v^{*}}
   R^{\tilde m_{\mathit{ls}} + c\,\mathit{ls}}
   \prod_{i\in\Hid} R_i^{\tilde m_i + c m_i}\\[2pt]
&= \underbrace{\Bigl(A'^{\,\tilde e} S^{\tilde v} R^{\tilde m_{\mathit{ls}}}
   \prod_{i\in\Hid} R_i^{\tilde m_i}\Bigr)}_{=\,T}
   \cdot
   \underbrace{\Bigl(A'^{\,e^{*}} S^{v^{*}} R^{\mathit{ls}}
   \prod_{i\in\Hid} R_i^{m_i}\Bigr)^{c}}_{=\,D^{c}\ \text{(Bổ đề \ref{lem:D})}} \\[2pt]
&= T\,D^{c}.
\end{align*}
Nhân hai vế với $D^{-c}$ được $\widehat T = T$. Do đó
$\Hash(\widehat T,A',n_v)=\Hash(T,A',n_v)=c$ và phép kiểm thành công.
\end{proof}

\begin{remark}[Vì sao không giả được]
Theo lập luận như Định lý~\ref{thm:soundness}, ai trả lời đúng cho hai thách thức khác
nhau trên cùng một $T$ sẽ để lộ một bộ $(e^{*},v^{*},\mathit{ls},\{m_i\})$ thoả đẳng thức
chữ ký --- tức là thật sự nắm một credential hợp lệ. Vì $c$ là giá trị băm của chính $T$,
kẻ gian phải cố định $T$ trước khi biết $c$, nên không thể dựng đáp số ngược.
\end{remark}

\section{Chống phát lại}

Số $n_v$ do verifier phát, dài $80$ bit, dùng đúng một lần và có thời hạn. Vì $n_v$ tham
gia vào $c=\Hash(T,A',n_v)$, một bằng chứng cũ gắn với $n_v$ cũ sẽ không khớp thách thức
mới. Tương tự, ở giai đoạn xin cấp, mỗi $\mathit{nonce}$ được issuer xoá ngay sau khi
dùng.

\section{Đối chiếu với mã nguồn}

\begin{center}
\begin{tabular}{@{}ll@{}}
\toprule
\textbf{Mục} & \textbf{Hàm trong mã nguồn} \\
\midrule
Sinh khoá & \texttt{issuer.setup} \\
Mã hoá thuộc tính & \texttt{common.encode\_attribute} \\
Hàm thách thức & \texttt{common.hash\_three} \\
Cam kết & \texttt{wallet.compute\_commitment} \\
Sigma --- bên chứng minh & \texttt{wallet.compute\_responses} \\
Sigma --- bên kiểm & \texttt{issuer.verify\_proof} \\
Ký mù & \texttt{issuer.sign\_blindly} \\
Gỡ mù và tự kiểm & \texttt{wallet.unblind\_signature}, \texttt{wallet.verify\_credential} \\
Bằng chứng trình bày & \texttt{wallet.create\_presentation} \\
Phép kiểm & \texttt{verifier.verify\_presentation} \\
\bottomrule
\end{tabular}
\end{center}

\section{Những gì lõi này chưa làm}

\begin{itemize}[leftmargin=1.2em,itemsep=4pt]
  \item \textbf{Predicate proof cho thuộc tính số.} Chưa chứng minh được ``đủ 18 tuổi''
        mà không lộ ngày sinh, vì ngày sinh hiện được băm nên mất thứ tự. Muốn làm cần mã
        hoá ngày sinh thành số có thứ tự kèm một range proof.
  \item \textbf{Revocation registry.} Chưa có cơ chế thu hồi thẻ đã cấp.
  \item \textbf{Ví chạy phía người dùng.} Trong bản demo, phần mật mã của ví chạy trên máy
        chủ. Đúng mô hình tự chủ danh tính thì phần này phải chạy trong trình duyệt để
        link secret không bao giờ rời thiết bị.
\end{itemize}

\end{document}
